\clearpage\fancyhead[L]{\small KOMET.}\hypertarget{native-378-lower}{}\label{native-378-lower}
\begin{center}\textbf{KOMET.}\par\vspace{4pt}\hrule\vspace{4pt}\fontsize{10}{12}\selectfont Göttingische gelehrte Anzeigen. Stück 19. S. 185\;..\;189. 31. Januar 1814.\par\vspace{4pt}\hrule\vspace{4pt}\end{center}

\par\indent Aus einem Brief vom 30. Dezember des vergangenen Jahres von Professor \textsc{Bessel} in Königsberg an Professor \textsc{Gauss} teilen wir hier einiges mit, das die astronomischen Leser unserer Zeitschrift gern früh erfahren werden. Aus verschiedenen Nachrichten ist bereits bekannt, dass die auch in unseren Anzeigen mehrmals angedeutete Aussicht (siehe 1811, Stück 130, und 1812, Stück 67 [S. 327 und 349 dieses Bandes]), der große Komet von 1811 könne im Sommer des Jahres mit vorzüglichen Instrumenten und unter günstigen Umständen noch einmal gesehen und beobachtet werden, nicht unbegründet war. Dieses Glück war den Astronomen an den europäischen Sternwarten trotz ihrer sorgfältigen Bemühungen nicht beschieden. Das lag teils an der zu hellen nächtlichen Dämmerung in dieser Jahreszeit, teils am weniger günstigen Klima, mit dem die meisten von ihnen zu kämpfen haben. Aus dem Kaukasus, aus \emph{Neu-Tscherkask}, erhalten die Astronomen die in ihrer Art einzigartigen Beobachtungen dieses bemerkenswerten Kometen. Herr von \textsc{Wisniwsky}, der seit mehreren Jahren in jenen abgelegenen Gegenden unterwegs ist, stellte sie fast ein Jahr nach dem Durchgang des Kometen durch die Sonnennähe an. Professor \textsc{Bessel}, der sich bekanntlich schon in mehr als einer Hinsicht um die Theorie dieses Kometen verdient gemacht und für ihn eine elliptische Bahn berechnet hatte, \SourceLeaf{379}erhielt diese wichtigen Beobachtungen aus Petersburg. Er wird sie später zur möglichst genauen Verfeinerung der Bahn verwenden. Die folgenden vorläufigen Ergebnisse verdienen es jedoch, schon jetzt bekannt gemacht zu werden.

\Needspace{14\baselineskip}\par\indent Herr von \textsc{Wisniwsky} beobachtete den Kometen vom 8. bis zum 27. August. In dieser Zeit erhielt er 29 Durchgänge durch das Kreismikrometer, dessen Feld er anhand mehrerer Sternpaare äußerst genau bestimmte. Er verglich den Kometen mit folgenden Sternen:

\begin{center}\fontsize{10}{12}\selectfont\setlength{\tabcolsep}{5pt}\renewcommand{\arraystretch}{1.15}\begin{tabular}{c|rr|rr}
& \multicolumn{2}{c|}{Rektaszension} & \multicolumn{2}{c}{Südl. Deklination}\\\hline
I & $324^\circ$ & $34'$ & $24^\circ$ & $31'$\\
II & 323 & 10 & 24 & 59\\
III & 322 & 37 & 25 & 16\\
IV & 322 & 45 & 25 & 57\\
V & 322 & 17 & 26 & 3\\
\end{tabular}\end{center}

\Needspace{15\baselineskip}\par\indent Nach der Reduktion von Professor \textsc{Bessel} ergeben sich folgende Positionen des Kometen relativ zu diesen Sternen:

\begin{center}\fontsize{9}{11}\selectfont\setlength{\tabcolsep}{3pt}\renewcommand{\arraystretch}{1.15}\begin{tabular}{l|c|c|c|rr|rr}
1812 & \shortstack{Mittlere Zeit in\\Neu-Tscherkask} & \shortstack{Vergl.\\Stern} & \shortstack{Anz. d.\\Beob.} & \multicolumn{2}{c|}{\shortstack{Rektaszension}} & \multicolumn{2}{c}{Deklination}\\\hline
Aug. 8 & $12^{\mathrm u}\;18^{\mathrm m}\;51^{\mathrm s}4$ & I & 4 & $-5'$ & $30''5$ & $-1'$ & $19''3$\\
11 & $12\;45\;21.7$ & II & 2 & $+25$ & 23.7 & $-3$ & 52.2\\
12 & $11\;42\;25.3$ & II & 7 & $+9$ & 7.5 & $-51$ & 15.3\\
12 & $12\;41\;41.0$ & III & 3 & $+41$ & 10.0 & $+1$ & 15.8\\
15 & $12\;10\;7.9$ & IV & 3 & $-18$ & 9.0 & $+11$ & 50.0\\
15 & $12\;30\;18.7$ & V & 5 & $+8$ & 39.8 & $+13$ & 0.0\\
17 & $13\;34\;8.4$ & V & 6 & $-25$ & 36.7 & $-1$ & 1.7\\
\end{tabular}\end{center}

\Needspace{12\baselineskip}\par\indent Von den Sternen selbst stehen nur drei (II, IV, V) in der \emph{Histoire céleste}. Professor \textsc{Bessel} reduzierte ihre dort angegebenen Positionen mit seinen eigenen Reduktionselementen auf das Jahr 1812 und fand:

\begin{center}\fontsize{10}{12}\selectfont\setlength{\tabcolsep}{4pt}\renewcommand{\arraystretch}{1.15}\begin{tabular}{c|rrr|rrr}
& \multicolumn{3}{c|}{Rektaszension} & \multicolumn{3}{c}{Südl. Deklination}\\\hline
II & $323^\circ$ & $9'$ & $1''5$ & $24^\circ$ & $59'$ & $26''5$\\
IV & 322 & 44 & 32.3 & 25 & 56 & 53.4\\
V & 322 & 16 & 17.1 & 26 & 3 & 36.9\\
\end{tabular}\end{center}

\Needspace{7\baselineskip}\par\indent Die Rektaszensionen von II und IV bestimmte er außerdem aus eigenen Beobachtungen:

\begin{center}\fontsize{10}{12}\selectfont\setlength{\tabcolsep}{4pt}\renewcommand{\arraystretch}{1.15}\begin{tabular}{lcrl}
II & $.\;.\;.\;.\;.\;.\;.\;.$ & $323^\circ\;9'\;5''3$ & (1 Beob.)\\
IV & $.\;.\;.\;.\;.\;.\;.\;.$ & $322\;44\;56.1$ & (2 Beob.)\\
\end{tabular}\end{center}

\Needspace{13\baselineskip}\par\indent Legt man bevorzugt diese Werte zugrunde, ergeben sich folgende beobachtete Positionen des Kometen:

\begin{center}\fontsize{9}{11}\selectfont\setlength{\tabcolsep}{3pt}\renewcommand{\arraystretch}{1.15}\begin{tabular}{l|c|rrr|rrr}
& Mittlere Zeit & \multicolumn{3}{c|}{Rektaszension} & \multicolumn{3}{c}{Südl. Deklination}\\\hline
Aug. 11 & $12^{\mathrm u}\;45^{\mathrm m}\;21^{\mathrm s}7$ & $323^\circ$ & $35'$ & $9''9$: & $25^\circ$ & $3'$ & $0''4$\\
12 & $11\;42\;25.3$ & 323 & 18 & 53.7 & 25 & 14 & 59.5\\
15 & $12\;10\;7.9$ & 322 & 27 & 7.9 & 25 & 44 & 45.3\\
15 & $12\;30\;18.7$ & 322 & 25 & 37.9 & 25 & 45 & 19.0\\
17 & $13\;34\;8.4$ & 321 & 51 & 21.4 & 26 & 4 & 20.7\\
\end{tabular}\end{center}

\Needspace{13\baselineskip}\par\indent Der Vergleich dieser Beobachtungen mit den rein elliptischen Elementen ergab für Professor \textsc{Bessel} folgende Differenzen
\SourceLeaf{380}\begin{center}\fontsize{10}{12}\selectfont\setlength{\tabcolsep}{8pt}\renewcommand{\arraystretch}{1.15}\begin{tabular}{l|c|c}
& \shortstack{Rektaszension} & Deklination\\\hline
Aug. 11 & $-31''2$ & $-10''0$:\\
12 & $-55.9$ & $+113.9$\\
15 & $-86.4$ & $+76.1$\\
15 & $-10.9$ & $+101.6$\\
17 & $-51.2$ & $+57.8$\\
\end{tabular}\end{center}

\par\indent Diese Übereinstimmung nach einem Abstand von mehr als drei Vierteljahren beweist, wie gut Herr \textsc{Bessel} die Beobachtungen zu nutzen wusste. Besonders diejenigen, die er selbst mit einem vorzüglichen Heliometer angestellt hatte, waren von ausgezeichneter Genauigkeit. Man sieht sogar, dass Herr \textsc{Bessel} selbst eine noch etwas bessere Übereinstimmung erzielt hätte, wenn er es schon damals der Mühe wert gefunden hätte, seine Elemente anhand der spätesten damaligen Beobachtungen noch einmal zu verfeinern, da diese bereits etwas in derselben Richtung von den Elementen abzuweichen begannen. Zugleich sieht man, dass die Störungen nur einen sehr kleinen Einfluss auf die Bewegung des Kometen gehabt haben können, mit dessen Berechnung Herr \textsc{Bessel} derzeit beschäftigt ist.
